\documentclass[12pt, a4paper]{article}

%% Pacotes
\usepackage[utf8]{inputenc}
\usepackage[T1]{fontenc}
\usepackage[brazil]{babel}
\usepackage[top=2.5cm, bottom=2.5cm, left=3cm, right=2.5cm]{geometry}
\usepackage{amsmath, amssymb, amsthm}
\usepackage{graphicx}
\usepackage{float}
\usepackage{caption}
\usepackage{subcaption}
\usepackage{booktabs}
\usepackage{multirow}
\usepackage{array}
\usepackage{xcolor}
\usepackage{hyperref}
\usepackage{setspace}
\usepackage{fancyhdr}
\usepackage{algorithm}
\usepackage{algpseudocode}
\usepackage{enumitem}
\usepackage{indentfirst}

%% Configurações gerais
\onehalfspacing
\setlength{\parindent}{1.25cm}

\hypersetup{
    colorlinks=true,
    linkcolor=black,
    citecolor=black,
    urlcolor=blue
}

\pagestyle{fancy}
\fancyhf{}
\rhead{\small Reconhecimento de Padrões Computacional -- UNESP}
\lhead{\small Atividade: Redução de Dimensionalidade, Classificação e Agrupamento}
\rfoot{\thepage}
\renewcommand{\headrulewidth}{0.4pt}

%% Macros matemáticas
\newcommand{\bx}{\mathbf{x}}
\newcommand{\bmu}{\boldsymbol{\mu}}
\newcommand{\bSigma}{\boldsymbol{\Sigma}}
\newcommand{\bW}{\mathbf{W}}
\newcommand{\bG}{\mathbf{G}}
\newcommand{\bY}{\mathbf{Y}}
\newcommand{\norm}[1]{\left\lVert#1\right\rVert}

%% Ambiente para resultados a preencher
\newcommand{\placeholder}[1]{\textcolor{gray}{\textit{#1}}}

% ===================================================================
\begin{document}

\begin{titlepage}
    \centering
    \vspace*{1.5cm}
    {\large \textbf{Universidade Estadual Paulista ``Júlio de Mesquita Filho''}} \\
    \vspace{0.3cm}
    {\large Instituto de Ciência e Tecnologia -- São José dos Campos} \\
    \vspace{0.3cm}
    {\large Programa de Pós-Graduação} \\[2cm]

    \rule{\textwidth}{1pt} \\[0.5cm]
    {\LARGE \textbf{Redução de Dimensionalidade, Classificação e Agrupamento}} \\[0.3cm]
    {\large Atividade Prática -- Reconhecimento de Padrões Computacional} \\[0.3cm]
    \rule{\textwidth}{1pt} \\[2.5cm]

    \begin{flushleft}
        \textbf{Aluno:} \underline{\hspace{8cm}} \\[0.5cm]
        \textbf{Matrícula:} \underline{\hspace{6cm}} \\[0.5cm]
        \textbf{Professor:} Prof.\ Dr.\ Rogério Galante Negri \\[0.5cm]
        \textbf{Disciplina:} Reconhecimento de Padrões Computacional (CLPR) \\[0.5cm]
        \textbf{Data:} \today
    \end{flushleft}

    \vfill
    {\large São José dos Campos -- SP, 2026}
\end{titlepage}

% ===================================================================
\tableofcontents
\newpage

% ===================================================================
\section{Introdução}
\label{sec:intro}

O Reconhecimento de Padrões (RP) é um campo interdisciplinar que visa extrair informação estruturada a partir de dados brutos~\cite{negri2021}. Em problemas do mundo real, os dados frequentemente apresentam alta dimensionalidade, o que impõe desafios tanto computacionais quanto estatísticos -- fenômeno amplamente conhecido como a \emph{maldição da dimensionalidade}~\cite{bishop2006}.

Esta atividade explora três grandes eixos do RP aplicados a dois conjuntos de dados de características distintas:

\begin{enumerate}[label=(\roman*)]
    \item \textbf{Redução de dimensionalidade:} empregando Análise de Componentes Principais (PCA, \emph{Principal Component Analysis}) e Seleção Sequencial Flutuante Progressiva (SFFS, \emph{Sequential Floating Forward Selection});
    \item \textbf{Classificação supervisionada:} com quatro métodos de complexidades e paradigmas diferentes;
    \item \textbf{Agrupamento não supervisionado:} com quatro métodos que cobrem abordagens sequenciais, fuzzy, probabilísticas e hierárquicas.
\end{enumerate}

O objetivo é analisar o impacto da redução dimensional no desempenho dos classificadores e comparar a capacidade de descoberta de estrutura latente nos dados pelos diferentes algoritmos de agrupamento.

% ===================================================================
\section{Conjuntos de Dados}
\label{sec:dados}

\subsection{Dados Artificiais (30 dimensões)}

O primeiro conjunto foi gerado com a função \texttt{make\_classification} do Scikit-learn com os seguintes parâmetros:

\begin{itemize}
    \item $N = 600$ amostras, $d = 30$ atributos, $c = 4$ classes balanceadas ($150$ amostras/classe);
    \item $10$ atributos informativos, $10$ redundantes (combinações lineares dos informativos), $5$ repetidos e $5$ puramente ruidosos;
    \item Semente aleatória: $42$.
\end{itemize}

A presença deliberada de redundância e ruído torna este conjunto adequado para avaliar técnicas de redução de dimensionalidade. A separação treino/avaliação seguiu a proporção $67\%/33\%$ (estratificada).

\subsection{Faces Olivetti}

O segundo conjunto é o \emph{Olivetti Faces Dataset}~\cite{samaria1994}, disponível no Scikit-learn:

\begin{itemize}
    \item $N = 400$ imagens, $d = 4096$ atributos (imagens $64 \times 64$ pixels em escala de cinza);
    \item $c = 40$ sujeitos distintos, $10$ imagens por sujeito (variações de iluminação, expressão e pose);
    \item Separação treino/avaliação: $67\%/33\%$ (estratificada).
\end{itemize}

A dimensionalidade extrema ($4096$) e o número elevado de classes ($40$) fazem deste conjunto um desafio representativo de aplicações reais de reconhecimento facial. Antes do processamento, os dados foram normalizados com \texttt{StandardScaler} ($\mu=0$, $\sigma^2=1$ por atributo).

% ===================================================================
\section{Redução de Dimensionalidade}
\label{sec:reducao}

\subsection{Análise de Componentes Principais (PCA)}
\label{sec:pca}

PCA é uma técnica linear não supervisionada que projeta os dados em um subespaço de menor dimensão, definido pelos autovetores da matriz de covariância associados aos maiores autovalores~\cite{bishop2006}. A projeção de uma amostra $\bx \in \mathbb{R}^d$ para o subespaço de $p$ componentes é:

\begin{equation}
    \bx' = \bW_p^T (\bx - \hat{\bmu}),
    \label{eq:pca}
\end{equation}

\noindent onde $\hat{\bmu}$ é a média empírica do conjunto de treino e $\bW_p = [\mathbf{w}_1, \ldots, \mathbf{w}_p] \in \mathbb{R}^{d \times p}$ contém os $p$ autovetores (componentes principais) ordenados de forma decrescente por autovalor.

O número de componentes $p$ é determinado automaticamente pelo critério de variância explicada acumulada:

\begin{equation}
    p = \min \left\{ k : \frac{\sum_{i=1}^{k} \lambda_i}{\sum_{i=1}^{d} \lambda_i} \geq 0{,}95 \right\},
    \label{eq:pca_var}
\end{equation}

\noindent onde $\lambda_i$ são os autovalores em ordem decrescente. O limiar de $95\%$ garante que a maior parte da variância seja preservada com redução substancial de dimensões.

\textbf{Importante:} $\hat{\bmu}$ e $\bW_p$ são estimados exclusivamente com os dados de treino e aplicados ao conjunto de avaliação, evitando vazamento de informação (\emph{data leakage}).

\subsection{Seleção Sequencial Flutuante Progressiva (SFFS)}
\label{sec:sffs}

SFFS é uma técnica supervisionada de seleção de atributos que alterna passos de adição (\emph{forward}) e remoção (\emph{backward}) de atributos em um subconjunto selecionado~\cite{theodoridis2008}. O critério de qualidade adotado é baseado nas matrizes de espalhamento inter e intraclasse:

\begin{equation}
    J(\Xi) = \frac{|\mathbf{V}_e + \mathbf{V}_i|}{|\mathbf{V}_i|},
    \label{eq:sffs_J}
\end{equation}

\noindent onde $\mathbf{V}_e$ é a matriz de espalhamento interclasse (\emph{between-class scatter}) e $\mathbf{V}_i$ é a matriz de espalhamento intraclasse (\emph{within-class scatter}):

\begin{align}
    \mathbf{V}_e &= \sum_{k=1}^{c} n_k (\hat{\bmu}_k - \hat{\bmu})(\hat{\bmu}_k - \hat{\bmu})^T, \label{eq:Ve} \\
    \mathbf{V}_i &= \sum_{k=1}^{c} \sum_{\bx \in \omega_k} (\bx - \hat{\bmu}_k)(\bx - \hat{\bmu}_k)^T. \label{eq:Vi}
\end{align}

Quanto maior $J(\Xi)$, maior a separabilidade entre classes em relação à compactação intraclasse. O Algoritmo~\ref{alg:sffs} descreve o procedimento completo.

\begin{algorithm}[H]
\caption{Sequential Floating Forward Selection (SFFS)}
\label{alg:sffs}
\begin{algorithmic}[1]
\Require Conjunto de atributos candidatos $\mathcal{F}$, número de atributos a selecionar $n_s$
\Ensure Subconjunto ótimo $\Xi^*$
\State $\Xi \leftarrow \emptyset$
\While{$|\Xi| < n_s$}
    \State \textbf{Passo Forward:} $\xi^+ \leftarrow \arg\max_{\xi \in \mathcal{F} \setminus \Xi} J(\Xi \cup \{\xi\})$
    \State $\Xi \leftarrow \Xi \cup \{\xi^+\}$
    \State \textbf{Passo Backward (floating):}
    \While{$|\Xi| > 1$}
        \State $\xi^- \leftarrow \arg\max_{\xi \in \Xi} J(\Xi \setminus \{\xi\})$
        \If{$J(\Xi \setminus \{\xi^-\}) > J(\Xi \setminus \{\xi^-\})_{\text{anterior}}$}
            \State $\Xi \leftarrow \Xi \setminus \{\xi^-\}$
        \Else
            \State \textbf{break}
        \EndIf
    \EndWhile
\EndWhile
\State \Return $\Xi$
\end{algorithmic}
\end{algorithm}

No caso do conjunto Olivetti Faces (4096 dimensões), o SFFS é aplicado sobre as componentes PCA (e não nos atributos originais), pois a avaliação do critério $J$ em espaço de 4096 dimensões é computacionalmente inviável.

\subsection{Resultados da Redução de Dimensionalidade}
\label{sec:res_reducao}

A Tabela~\ref{tab:reducao} apresenta o número de dimensões antes e após cada método de redução.

\begin{table}[H]
    \centering
    \caption{Dimensões resultantes após PCA e SFFS.}
    \label{tab:reducao}
    \begin{tabular}{lccc}
        \toprule
        \textbf{Dataset} & \textbf{Original} & \textbf{PCA (95\%)} & \textbf{SFFS} \\
        \midrule
        Artificial       & 30     & \placeholder{--}  & \placeholder{--} \\
        Olivetti Faces   & 4096   & \placeholder{--}  & \placeholder{--} \\
        \bottomrule
    \end{tabular}
\end{table}

\begin{figure}[H]
    \centering
    \fbox{\parbox{0.45\textwidth}{\centering\vspace{3cm}
        \textit{Variância explicada acumulada -- Artificial}
        \vspace{3cm}}}
    \hfill
    \fbox{\parbox{0.45\textwidth}{\centering\vspace{3cm}
        \textit{Variância explicada acumulada -- Olivetti}
        \vspace{3cm}}}
    \caption{Variância explicada acumulada pelas componentes PCA. A linha tracejada indica o limiar de 95\%.}
    \label{fig:pca_var}
\end{figure}

% ===================================================================
\section{Classificação Supervisionada}
\label{sec:classificacao}

Quatro métodos de classificação foram aplicados a cada configuração de dados (original, PCA e SFFS). Os classificadores foram treinados exclusivamente com os dados de treino e avaliados no conjunto de avaliação.

\subsection{Rede RBF com Mapeamento Polinomial e Saída SSE (Atividade III)}
\label{sec:polyrbf}

Este classificador combina três estágios em sequência~\cite{negri2021}:

\textbf{Estágio 1 -- Mapeamento Polinomial:} o vetor de entrada $\bx \in \mathbb{R}^d$ é expandido para o espaço de características polinomiais de grau $q=2$:
\begin{equation}
    \bx \mapsto \boldsymbol{\phi}(\bx) = [1,\; x_1,\; x_2,\; \ldots,\; x_d,\; x_1^2,\; x_1 x_2,\; \ldots,\; x_d^2]^T \in \mathbb{R}^{d'},
    \label{eq:poly}
\end{equation}
com $d' = \binom{d+q}{q}$. Para $d=30$ e $q=2$, obtém-se $d' = 496$ features expandidas.

\textbf{Estágio 2 -- Camada Oculta RBF:} os centróides $\{\bmu_j\}_{j=1}^{K}$ são estimados via K-Means no espaço polinomial. A ativação de cada neurônio oculto é:
\begin{equation}
    g_j(\bx) = \exp\!\left(-\frac{\norm{\boldsymbol{\phi}(\bx) - \bmu_j}^2}{2\sigma_j^2}\right), \quad j = 1, \ldots, K,
    \label{eq:rbf_act}
\end{equation}
com $\sigma_j^2$ estimado como a distância média ao centróide mais próximo.

\textbf{Estágio 3 -- Saída SSE Linear:} os pesos da camada de saída são obtidos analiticamente pela solução de mínimos quadrados~\cite{negri2021}:
\begin{equation}
    \bW = (\bG^T \bG)^{-1} \bG^T \bY,
    \label{eq:sse}
\end{equation}
onde $\bG \in \mathbb{R}^{N \times K}$ é a matriz de ativações da camada oculta e $\bY \in \mathbb{R}^{N \times c}$ é a codificação \emph{one-hot} dos rótulos. A classe predita é $\hat{y} = \arg\max_k [\bG \bW]_k$.

\subsection{Neurônios Aleatórios -- ELM (Atividade IV)}
\label{sec:elm}

O método de Neurônios Aleatórios (Extreme Learning Machine, ELM) é uma variante da Rede RBF onde os centróides $\bmu_j$ são amostrados aleatoriamente do conjunto de treino, sem a fase de K-Means~\cite{negri2021}. Os parâmetros da camada oculta são \emph{fixos}, e apenas os pesos de saída $\bW$ são aprendidos, também pela Equação~\eqref{eq:sse}:
\begin{equation}
    \bmu_j \sim \text{Uniforme}(\{\bx_1, \ldots, \bx_N\}), \quad j = 1, \ldots, K.
    \label{eq:elm_centers}
\end{equation}

O desvio padrão da função de ativação é estimado por:
\begin{equation}
    \sigma = \hat{\sigma}(\bX) \cdot \sqrt{\frac{d}{K}},
    \label{eq:elm_sigma}
\end{equation}
onde $\hat{\sigma}(\bX)$ é o desvio padrão global dos dados de treino e $K$ é o número de neurônios ocultos.

A principal vantagem do ELM sobre a RBFNet é o custo computacional reduzido pela eliminação do K-Means. A principal desvantagem é a maior variância dos resultados, pois os centróides não são otimizados para a distribuição dos dados.

\subsection{SVM com Kernel RBF}
\label{sec:svm}

A SVM (\emph{Support Vector Machine}) busca o hiperplano de máxima margem no espaço de características induzido pelo kernel~\cite{bishop2006}. O kernel Gaussiano (RBF) é definido por:
\begin{equation}
    K(\bx_i, \bx_j) = \exp(-\gamma \norm{\bx_i - \bx_j}^2),
    \label{eq:kernel_rbf}
\end{equation}
com $\gamma = $ \texttt{scale} (i.e., $\gamma = 1/(d \cdot \text{Var}(\bX))$) e penalidade $C = 10$. Para problemas multiclasse ($c > 2$), adota-se a estratégia \emph{One-vs-Rest} (OVR), treinando $c$ SVMs binárias.

\subsection{K-Vizinhos Mais Próximos (KNN)}
\label{sec:knn}

O KNN é um classificador não-paramétrico e baseado em instâncias. A classe de uma amostra $\bx$ é determinada por votação majoritária entre seus $k$ vizinhos mais próximos (distância Euclidiana) no conjunto de treino:
\begin{equation}
    \hat{y}(\bx) = \arg\max_{\omega_j} \sum_{i \in \mathcal{N}_k(\bx)} \mathbf{1}[y_i = \omega_j],
    \label{eq:knn}
\end{equation}
onde $\mathcal{N}_k(\bx)$ é o conjunto dos $k$ vizinhos mais próximos. Nesta atividade, $k = 5$.

\subsection{Métricas de Avaliação para Classificação}
\label{sec:metricas_class}

\textbf{Acerto Global (OA):}
\begin{equation}
    A_g = \frac{1}{m} \sum_{i=1}^{c} a_{ii},
    \label{eq:oa}
\end{equation}
onde $m$ é o total de amostras avaliadas e $a_{ii}$ são os elementos da diagonal principal da matriz de confusão.

\textbf{Coeficiente Kappa:}
\begin{equation}
    \kappa = \frac{A_g - A_a}{1 - A_a}, \quad A_a = \frac{1}{m^2} \sum_{i=1}^{c} \left(\sum_j a_{ij}\right)\!\left(\sum_j a_{ji}\right),
    \label{eq:kappa}
\end{equation}
onde $A_a$ representa o acerto esperado por acaso. $\kappa \in [0,1]$ com $\kappa = 1$ indicando concordância perfeita e $\kappa = 0$ indicando concordância aleatória~\cite{negri2021}.

\subsection{Resultados de Classificação}
\label{sec:res_class}

As Tabelas~\ref{tab:class_artificial} e \ref{tab:class_olivetti} apresentam os resultados (OA e $\kappa$) para os quatro classificadores, nas três configurações de espaço de atributos.

\begin{table}[H]
    \centering
    \caption{Classificação -- Dataset Artificial (30D). OA: Acerto Global; $\kappa$: Coeficiente Kappa.}
    \label{tab:class_artificial}
    \setlength{\tabcolsep}{5pt}
    \begin{tabular}{lcccccc}
        \toprule
        \multirow{2}{*}{\textbf{Método}} &
        \multicolumn{2}{c}{\textbf{Original (30D)}} &
        \multicolumn{2}{c}{\textbf{PCA}} &
        \multicolumn{2}{c}{\textbf{SFFS}} \\
        \cmidrule(lr){2-3} \cmidrule(lr){4-5} \cmidrule(lr){6-7}
        & OA & $\kappa$ & OA & $\kappa$ & OA & $\kappa$ \\
        \midrule
        PolyMap-RBFNet-SSE & \placeholder{--} & \placeholder{--} & \placeholder{--} & \placeholder{--} & \placeholder{--} & \placeholder{--} \\
        Random Neurons (ELM) & \placeholder{--} & \placeholder{--} & \placeholder{--} & \placeholder{--} & \placeholder{--} & \placeholder{--} \\
        SVM-RBF ($C=10$)  & \placeholder{--} & \placeholder{--} & \placeholder{--} & \placeholder{--} & \placeholder{--} & \placeholder{--} \\
        KNN ($k=5$)        & \placeholder{--} & \placeholder{--} & \placeholder{--} & \placeholder{--} & \placeholder{--} & \placeholder{--} \\
        \bottomrule
    \end{tabular}
\end{table}

\begin{table}[H]
    \centering
    \caption{Classificação -- Olivetti Faces (4096D). OA: Acerto Global; $\kappa$: Coeficiente Kappa.}
    \label{tab:class_olivetti}
    \setlength{\tabcolsep}{5pt}
    \begin{tabular}{lcccccc}
        \toprule
        \multirow{2}{*}{\textbf{Método}} &
        \multicolumn{2}{c}{\textbf{Original}} &
        \multicolumn{2}{c}{\textbf{PCA}} &
        \multicolumn{2}{c}{\textbf{SFFS (sobre PCA)}} \\
        \cmidrule(lr){2-3} \cmidrule(lr){4-5} \cmidrule(lr){6-7}
        & OA & $\kappa$ & OA & $\kappa$ & OA & $\kappa$ \\
        \midrule
        PolyMap-RBFNet-SSE & \placeholder{--} & \placeholder{--} & \placeholder{--} & \placeholder{--} & \placeholder{--} & \placeholder{--} \\
        Random Neurons (ELM) & \placeholder{--} & \placeholder{--} & \placeholder{--} & \placeholder{--} & \placeholder{--} & \placeholder{--} \\
        SVM-RBF ($C=10$)  & \placeholder{--} & \placeholder{--} & \placeholder{--} & \placeholder{--} & \placeholder{--} & \placeholder{--} \\
        KNN ($k=5$)        & \placeholder{--} & \placeholder{--} & \placeholder{--} & \placeholder{--} & \placeholder{--} & \placeholder{--} \\
        \bottomrule
    \end{tabular}
\end{table}

\begin{figure}[H]
    \centering
    \fbox{\parbox{0.45\textwidth}{\centering\vspace{3cm}
        \textit{Matriz de confusão -- melhor classificador, Dataset Artificial}
        \vspace{3cm}}}
    \hfill
    \fbox{\parbox{0.45\textwidth}{\centering\vspace{3cm}
        \textit{Matriz de confusão -- melhor classificador, Olivetti}
        \vspace{3cm}}}
    \caption{Matrizes de confusão do melhor classificador em cada dataset.}
    \label{fig:matconf}
\end{figure}

% ===================================================================
\section{Agrupamento Não Supervisionado}
\label{sec:agrupamento}

O agrupamento foi aplicado aos dados sem utilizar os rótulos durante o treinamento. Quatro métodos foram comparados, dois deles obrigatórios da Atividade V.

\subsection{BSAS -- Basic Sequential Algorithmic Scheme (Atividade V)}
\label{sec:bsas}

O BSAS é um algoritmo de agrupamento sequencial que processa cada amostra uma única vez~\cite{theodoridis2008}. Um novo cluster é criado sempre que a distância mínima a todos os clusters existentes supera um limiar $\Theta$:

\begin{algorithm}[H]
\caption{Basic Sequential Algorithmic Scheme (BSAS)}
\label{alg:bsas}
\begin{algorithmic}[1]
\Require Dados $\{\bx_1, \ldots, \bx_N\}$, limiar $\Theta$, máximo de clusters $c_{\max}$
\State $m \leftarrow 1$; $C_1 \leftarrow \{\bx_1\}$; $\bmu_1 \leftarrow \bx_1$
\For{$i = 2$ to $N$}
    \State $k^* \leftarrow \arg\min_{k \in \{1,\ldots,m\}} d(\bx_i, \bmu_k)$
    \If{$d(\bx_i, \bmu_{k^*}) > \Theta$ \textbf{and} $m < c_{\max}$}
        \State $m \leftarrow m + 1$; $C_m \leftarrow \{\bx_i\}$; $\bmu_m \leftarrow \bx_i$
    \Else
        \State $C_{k^*} \leftarrow C_{k^*} \cup \{\bx_i\}$
        \State $\bmu_{k^*} \leftarrow \frac{(n_{k^*}-1)\,\bmu_{k^*} + \bx_i}{n_{k^*}}$ \Comment{Atualização incremental}
    \EndIf
\EndFor
\end{algorithmic}
\end{algorithm}

A atualização incremental do centróide mantém a complexidade linear em memória. O BSAS é sensível à ordem dos dados e ao valor de $\Theta$, que deve ser ajustado com base nas distâncias do conjunto.

\subsection{Fuzzy K-Means -- FKM (Atividade V)}
\label{sec:fkm}

O FKM (ou FCM -- \emph{Fuzzy C-Means}) estende o K-Means para graus de pertinência contínuos~\cite{theodoridis2008}. A função objetivo é:
\begin{equation}
    J_{\text{FKM}} = \sum_{i=1}^{N} \sum_{k=1}^{c} \lambda_{ik}^\beta \norm{\bx_i - \bmu_k}^2,
    \label{eq:fkm_obj}
\end{equation}
onde $\lambda_{ik} \in [0,1]$ é o grau de pertinência da amostra $i$ ao cluster $k$, e $\beta > 1$ é o parâmetro de fuzzificação (aqui, $\beta = 2$).

As atualizações iterativas são:
\begin{align}
    \lambda_{ik} &= \left[\sum_{j=1}^{c} \left(\frac{\norm{\bx_i - \bmu_k}}{\norm{\bx_i - \bmu_j}}\right)^{\!\frac{2}{\beta-1}}\right]^{-1}, \label{eq:fkm_lambda} \\
    \bmu_k &= \frac{\sum_{i=1}^{N} \lambda_{ik}^\beta \bx_i}{\sum_{i=1}^{N} \lambda_{ik}^\beta}. \label{eq:fkm_mu}
\end{align}

O algoritmo converge quando $\max_{i,k} |\lambda_{ik}^{(t+1)} - \lambda_{ik}^{(t)}| < \varepsilon$, com $\varepsilon = 10^{-4}$.

\subsection{K-Means}
\label{sec:kmeans}

O K-Means minimiza a variância intraclasse:
\begin{equation}
    J_{\text{KM}} = \sum_{k=1}^{c} \sum_{\bx \in C_k} \norm{\bx - \bmu_k}^2.
    \label{eq:kmeans_obj}
\end{equation}

Utilizou-se a implementação do Scikit-learn com $n\_\text{init} = 10$ reinicializações aleatórias e critério \texttt{k-means++} para inicialização dos centróides.

\subsection{Agrupamento Hierárquico -- Ward}
\label{sec:ward}

O agrupamento hierárquico aglomerativo com ligação de Ward funde iterativamente os dois clusters cuja fusão minimiza o aumento da variância intraclasse~\cite{bishop2006}:
\begin{equation}
    \Delta(C_a, C_b) = \frac{n_a n_b}{n_a + n_b} \norm{\bmu_a - \bmu_b}^2,
    \label{eq:ward}
\end{equation}
onde $n_a$, $n_b$ são os tamanhos e $\bmu_a$, $\bmu_b$ os centróides dos clusters mesclados.

\subsection{Métricas de Avaliação para Agrupamento}
\label{sec:metricas_cluster}

\textbf{Silhouette Coefficient (SC):}
\begin{equation}
    s(i) = \frac{b(i) - a(i)}{\max\{a(i),\, b(i)\}},
    \label{eq:silhouette}
\end{equation}
onde $a(i)$ é a distância média intracluster e $b(i)$ é a menor distância média para os demais clusters. $SC = (1/N)\sum_i s(i) \in [-1, 1]$.

\textbf{Índice Calinski-Harabász (CHI):}
\begin{equation}
    \text{CHI} = \frac{\text{tr}(\mathbf{V}_e)}{\text{tr}(\mathbf{V}_i)} \cdot \frac{N - c}{c - 1}.
    \label{eq:chi}
\end{equation}
Valores maiores indicam clusters mais compactos e bem separados.

\textbf{V-measure} (requer rótulos de referência):
\begin{equation}
    V = \frac{2 \cdot H_c \cdot C_c}{H_c + C_c},
    \label{eq:vmeasure}
\end{equation}
onde $H_c$ é a homogeneidade (cada cluster contém apenas membros de uma classe) e $C_c$ é a completude (todos os membros de uma classe estão no mesmo cluster).

\subsection{Resultados de Agrupamento}
\label{sec:res_cluster}

\begin{table}[H]
    \centering
    \caption{Agrupamento -- Dataset Artificial ($c=4$ clusters). SC: Silhouette; CHI: Calinski-Harabász; V: V-measure.}
    \label{tab:cluster_artificial}
    \begin{tabular}{lcccc}
        \toprule
        \textbf{Método} & \textbf{$k$} & \textbf{SC} & \textbf{CHI} & \textbf{V-measure} \\
        \midrule
        BSAS             & \placeholder{--} & \placeholder{--} & \placeholder{--} & \placeholder{--} \\
        Fuzzy K-Means    & 4                & \placeholder{--} & \placeholder{--} & \placeholder{--} \\
        K-Means          & 4                & \placeholder{--} & \placeholder{--} & \placeholder{--} \\
        Ward Hierárquico & 4                & \placeholder{--} & \placeholder{--} & \placeholder{--} \\
        \bottomrule
    \end{tabular}
\end{table}

\begin{table}[H]
    \centering
    \caption{Agrupamento -- Olivetti Faces ($c=10$ clusters). SC: Silhouette; CHI: Calinski-Harabász; V: V-measure (em relação aos 40 rótulos verdadeiros).}
    \label{tab:cluster_olivetti}
    \begin{tabular}{lcccc}
        \toprule
        \textbf{Método} & \textbf{$k$} & \textbf{SC} & \textbf{CHI} & \textbf{V-measure} \\
        \midrule
        BSAS             & \placeholder{--} & \placeholder{--} & \placeholder{--} & \placeholder{--} \\
        Fuzzy K-Means    & 10               & \placeholder{--} & \placeholder{--} & \placeholder{--} \\
        K-Means          & 10               & \placeholder{--} & \placeholder{--} & \placeholder{--} \\
        Ward Hierárquico & 10               & \placeholder{--} & \placeholder{--} & \placeholder{--} \\
        \bottomrule
    \end{tabular}
\end{table}

\begin{figure}[H]
    \centering
    \fbox{\parbox{0.45\textwidth}{\centering\vspace{3cm}
        \textit{Dendrograma (Ward) -- Dataset Artificial}
        \vspace{3cm}}}
    \hfill
    \fbox{\parbox{0.45\textwidth}{\centering\vspace{3cm}
        \textit{Projeção 2D PCA com rótulos de agrupamento -- Olivetti}
        \vspace{3cm}}}
    \caption{Visualizações dos resultados de agrupamento.}
    \label{fig:cluster_vis}
\end{figure}

% ===================================================================
\section{Discussão}
\label{sec:discussao}

\subsection{Impacto da Maldição da Dimensionalidade}

O dataset Olivetti Faces, com $d = 4096$ atributos e apenas $N = 400$ amostras, representa um regime em que $d \gg N$. Neste regime, a distância Euclidiana perde poder discriminativo~\cite{bishop2006}: a razão $d_{\max}/d_{\min}$ (distância máxima sobre mínima entre pares) tende a $1$ à medida que $d \to \infty$, tornando o espaço praticamente equidistante. Isso afeta especialmente o KNN e o BSAS, que dependem fortemente de distâncias no espaço original.

O PCA mitiga este problema ao projetar os dados em um subespaço de menor dimensão que concentra a maior parte da variância. \placeholder{Discutir aqui, com base nos resultados obtidos, o ganho de desempenho após PCA.}

\subsection{PCA \textit{vs.}\ SFFS: Extração \textit{vs.}\ Seleção}

PCA é uma técnica de \emph{extração} de características: as novas variáveis (componentes principais) são combinações lineares de todas as originais, sem interpretabilidade direta. Por outro lado, o SFFS realiza \emph{seleção} de atributos: retém um subconjunto das variáveis originais (ou das componentes PCA, neste caso), preservando a interpretabilidade.

A vantagem do SFFS está em seu critério supervisionado: o critério $J$ (Equação~\ref{eq:sffs_J}) avalia explicitamente a discriminabilidade entre classes, enquanto o PCA é insensível aos rótulos. Isso pode fazer o SFFS encontrar subconjuntos mais relevantes para classificação em cenários específicos.

\placeholder{Discutir aqui os atributos/componentes selecionados pelo SFFS em cada dataset e compará-los com os componentes PCA retidos.}

\subsection{Classificação: RBFNet \textit{vs.}\ ELM}

A diferença fundamental entre o PolyMap-RBFNet-SSE e o Random Neurons (ELM) é a forma de determinar os centróides da camada oculta: K-Means (supervisionado no posicionamento) \textit{vs.}\ amostras aleatórias do conjunto de treino. Em geral, espera-se que o K-Means produza centróides mais representativos da distribuição dos dados, mas o ELM tem vantagem computacional e pode ser suficiente quando $K$ é suficientemente grande~\cite{negri2021}.

\placeholder{Comparar os resultados de OA e $\kappa$ dos dois métodos nos dois datasets e discutir a relação custo-benefício.}

\subsection{Agrupamento: Métodos Hard \textit{vs.}\ Fuzzy}

O FKM, ao contrário do K-Means, permite que amostras pertençam parcialmente a mais de um cluster. Isso é especialmente relevante em dados com sobreposição entre classes -- como o dataset Artificial, que contém atributos redundantes e ruidosos. O parâmetro $\beta=2$ controla o grau de fuzzificação: $\beta \to 1$ recupera o K-Means rígido, e $\beta \to \infty$ distribui uniformemente as pertinências.

\placeholder{Discutir aqui se o FKM apresentou ganhos em relação ao K-Means nos datasets analisados, com base nos valores de SC, CHI e V-measure.}

\subsection{BSAS: Sensibilidade ao Limiar $\Theta$}

O BSAS é o único método analisado que determina automaticamente o número de clusters (limitado a $c_{\max}$), dado o limiar $\Theta$. Isso é uma vantagem em cenários onde $c$ é desconhecido, mas exige calibração cuidadosa: valores de $\Theta$ muito baixos fragmentam clusters verdadeiros, enquanto valores muito altos os fundem. \placeholder{Descrever o valor de $\Theta$ adotado e como ele foi determinado.}

% ===================================================================
\section{Conclusão}
\label{sec:conclusao}

Esta atividade investigou o efeito da redução de dimensionalidade sobre o desempenho de quatro classificadores e avaliou a capacidade de quatro algoritmos de agrupamento em recuperar estrutura latente em dois datasets com características distintas.

Os resultados confirmam que a alta dimensionalidade -- especialmente no dataset Olivetti Faces -- apresenta desafios que métodos como PCA e SFFS ajudam a mitigar. \placeholder{Completar com as principais conclusões derivadas dos resultados obtidos: qual método de classificação foi superior, qual configuração dimensional foi mais adequada, e qual algoritmo de agrupamento revelou melhor estrutura.}

A comparação entre os métodos obrigatórios da Atividade III (PolyMap-RBFNet-SSE) e da Atividade IV (Random Neurons -- ELM) evidencia o trade-off entre custo computacional e precisão. \placeholder{Concluir sobre o mérito relativo de cada abordagem no contexto dos datasets analisados.}

Para trabalhos futuros, seria interessante explorar: (i) outros critérios de parada para o SFFS; (ii) outras variantes do FKM com $\beta > 2$; (iii) avaliação do impacto do número de neurônios ocultos $K$ no desempenho do ELM e da RBFNet.

% ===================================================================
\bibliographystyle{plain}
\begin{thebibliography}{9}

\bibitem{negri2021}
Negri, R.\ G.
\newblock \textit{Reconhecimento de Padrões Computacional}.
\newblock Blucher, São Paulo, 2021.

\bibitem{bishop2006}
Bishop, C.\ M.
\newblock \textit{Pattern Recognition and Machine Learning}.
\newblock Springer, New York, 2006.

\bibitem{theodoridis2008}
Theodoridis, S.; Koutroumbas, K.
\newblock \textit{Pattern Recognition}.
\newblock 4th ed. Academic Press, Burlington, 2008.

\bibitem{samaria1994}
Samaria, F.; Harter, A.
\newblock Parameterisation of a stochastic model for human face identification.
\newblock In: \textit{Proceedings of 2nd IEEE Workshop on Applications of Computer Vision}, pp.\ 138--142, 1994.

\end{thebibliography}

\end{document}
